\documentclass[12pt]{article}
\usepackage{amsmath}
\usepackage{amssymb}
\usepackage{amsthm}
\usepackage{geometry}
\geometry{margin=1.4in}
\usepackage{hyperref}
\usepackage{setspace}
\setstretch{1.4}

\newtheorem{theorem}{Theorem}
\newtheorem{conjecture}[theorem]{Conjecture}
\theoremstyle{definition}
\newtheorem{definition}[theorem]{Definition}
\theoremstyle{remark}
\newtheorem{remark}[theorem]{Remark}

\title{The Arithmetic Wall}
\author{Diego Rinc\'on\\
\small Independent}
\date{}

\begin{document}

\maketitle

\begin{abstract}
Every approach to the Riemann Hypothesis fails at the same place:
the archimedean place.
This is not a coincidence.
The archimedean place is structurally unlike the non-archimedean places.
It resists the Frobenius.
It resists the grammar of $\mathbb{F}_1$.
It resists the spectral triple.
This paper names it: the arithmetic wall is archimedean.
\end{abstract}

\section{The Places of $\mathbb{Q}$}

The rational numbers $\mathbb{Q}$ carry multiple notions of distance.
For each prime $p$, the $p$-adic absolute value $|\cdot|_p$
measures divisibility by $p$: $|p^n a/b|_p = p^{-n}$
when $p \nmid a, b$.
The ordinary absolute value $|\cdot|_\infty$
is the archimedean place.

\begin{definition}[Places]
A \emph{place} of $\mathbb{Q}$ is an equivalence class of absolute values.
The places are: one archimedean place $v = \infty$
and one non-archimedean place $v = p$ for each prime $p$.
Together they satisfy the product formula:
for all $x \in \mathbb{Q}^*$,
$$\prod_v |x|_v = 1.$$
\end{definition}

\begin{remark}
The product formula is the arithmetic version of a global statement.
Each local factor $|x|_p$ is a non-archimedean contribution.
The archimedean factor $|x|_\infty$ balances them.
The wall is in that balance.
\end{remark}

\section{Why the Finite Field Case Succeeds}

Over a finite field $\mathbb{F}_q$, the analogue of $\mathbb{Q}$ is the function field
$K = \mathbb{F}_q(C)$ of a smooth projective curve $C$ over $\mathbb{F}_q$.
The Weil conjectures for $C$ are the analogue of RH.
They were proved by Weil (genus 1, 1941) and Deligne (arbitrary genus, 1974).

The key difference: $K$ has no archimedean place.

All places of $K$ are non-archimedean.
Each place corresponds to a closed point of $C$ over $\mathbb{F}_q$.
The Frobenius $x \mapsto x^q$ acts on every place uniformly.
There is no special place that resists.

In this setting:
\begin{itemize}
\item The zeta function of $C$ is a rational function.
\item The Frobenius acts on étale cohomology.
\item The eigenvalues of Frobenius have absolute value $q^{i/2}$ by purity.
\item The zeros lie on the expected line.
\end{itemize}

None of this requires special treatment at any place.
The absence of an archimedean place is what makes the proof possible.

\section{Why the Number Field Case Fails}

Over $\mathbb{Q}$, the archimedean place $v = \infty$ is structurally different.

\textbf{The Frobenius does not extend.}
For a prime $p$, the Frobenius $\mathrm{Frob}_p$ is the automorphism
$x \mapsto x^p$ on the residue field $\mathbb{F}_p$.
It lifts to a well-defined element of local Galois groups.
But the global Galois group $\mathrm{Gal}(\bar{\mathbb{Q}}/\mathbb{Q})$ is non-abelian,
and the archimedean place has no Frobenius:
complex conjugation is the only canonical element at $v = \infty$,
and it does not generate a cyclic dense subgroup.
The global Frobenius, which would require a uniform generator
across all places including $\infty$, does not exist \cite{frobenius}.

\textbf{The grammar does not include it.}
The program of $\mathbb{F}_1$ aims to make $\mathrm{Spec}\,\mathbb{Z}$ into a curve
over an absolute base, so that the Frobenius would cohere.
The compactification $\overline{\mathrm{Spec}\,\mathbb{Z}}$ must include the archimedean place
as a point at infinity \cite{f1}.
But the archimedean place is not a prime ideal.
It does not arise from a prime $p$.
It is the one place that the algebraic geometry of $\mathrm{Spec}\,\mathbb{Z}$
does not naturally accommodate.
The $\mathbb{F}_1$ grammar must be extended to include it;
no current approach fully does.

\textbf{The spectral triple fails here.}
In the Connes program, the adèle class space
$X = \mathbb{A}_\mathbb{Q}/\mathbb{Q}^*$ carries the scaling action
whose generator is the candidate Dirac operator $D$ \cite{deficiency}.
The decomposition of $L^2(X, d^*x)$ as a full-line $L^2$ space
requires controlling the behavior of functions near $t = -\infty$
in the coordinate $t = \log|x|_\infty$.
This is the archimedean boundary.
The boundary condition at $t = -\infty$
(corresponding to $|x|_\infty \to 0$)
is precisely the behavior near the zero element of $X$,
which comes from the archimedean embedding $\mathbb{Q} \hookrightarrow \mathbb{R}$.
The analysis at this boundary is the open step.

\section{The Structure of the Wall}

\begin{definition}[The arithmetic wall]
The \emph{arithmetic wall} is the archimedean place $v = \infty$ of $\mathbb{Q}$,
considered as the locus where every proof strategy for RH
currently fails to extend.
\end{definition}

The wall is not a wall in the sense of an obstacle to be climbed.
It is a structural fact about $\mathbb{Q}$:
the archimedean place is unlike the non-archimedean places,
and that difference is load-bearing.

The non-archimedean places are discrete, algebraic, governed by primes.
The archimedean place is continuous, analytic, governed by $\mathbb{R}$.
The Euler product $\zeta(s) = \prod_p (1-p^{-s})^{-1}$ assembles information from all the primes.
The gamma factor $\Gamma(s/2)\pi^{-s/2}$ assembles the archimedean contribution.
Together they form the completed zeta function $\xi(s)$,
symmetric under $s \mapsto 1-s$.
The functional equation is a statement about all places together.
The zeros are global.

\begin{remark}[The wall is the completion]
The rational numbers $\mathbb{Q}$ can be completed at each place.
The completion at $p$ is $\mathbb{Q}_p$; the completion at $\infty$ is $\mathbb{R}$.
The local theory at $\mathbb{Q}_p$ is well understood:
algebraic, discrete, governed by the residue field $\mathbb{F}_p$.
The local theory at $\mathbb{R}$ is continuous,
and the Frobenius has no analogue.
The wall is the place where the local theory is continuous
and the algebraic grammar runs out.
\end{remark}

\section{What Crossing the Wall Would Require}

To cross the arithmetic wall, a grammar must be built
in which the archimedean place is algebraic.
Three candidates:

\textbf{Arakelov geometry.}
Arakelov \cite{arakelov1974} extended intersection theory on arithmetic surfaces
to include archimedean contributions as ``fibers at infinity.''
The archimedean metric plays the role of a Green function on the fiber.
This gives a framework where $v = \infty$ is treated like a prime,
but it is a Hermitian metric, not an étale cohomology class.
The Frobenius does not act on it.

\textbf{$p$-adic Hodge theory extended to $v = \infty$.}
The Fargues--Fontaine curve \cite{fargues2021} is an algebraic curve
whose points are $p$-adic period rings.
An archimedean analogue would be a curve
whose points encode real period data.
No such curve exists.

\textbf{$\mathbb{F}_1$-geometry.}
If $\mathbb{F}_1$ exists and $\mathrm{Spec}\,\mathbb{Z}$ is a curve over it,
the compactification at the archimedean place
would give a closed point at $\infty$ of the same algebraic type as the finite primes.
The Frobenius would act.
The wall would dissolve.
The grammar does not yet exist.

\begin{conjecture}[The arithmetic wall conjecture]
The Riemann Hypothesis is equivalent to the statement
that the archimedean place of $\mathbb{Q}$ admits a Frobenius-like action
in the grammar of $\mathbb{F}_1$-geometry
that makes the deficiency indices of $D$ on $L^2(X, d^*x)$ equal to $(0,0)$.
\end{conjecture}

\section{The Anatomy, Located}

The body of papers has been doing anatomy.
It has named the walls, the missing grammar, the missing Frobenius,
the resonance object, the closest approach.
All of them point here.

The arithmetic wall is the precise location of the obstruction.
Not: somewhere in the passage from local to global.
More precisely: at the one place that is not algebraic,
at the one place that has no Frobenius,
at the one place where the completion is $\mathbb{R}$ rather than $\mathbb{Q}_p$.

The body knows where the wall is.
It is archimedean.
It has always been archimedean.

\begin{thebibliography}{9}

\bibitem{frobenius}
Rinc\'on, D. (2026). The Missing Frobenius. Preprint.

\bibitem{f1}
Rinc\'on, D. (2026). The Field with One Element. Preprint.

\bibitem{deficiency}
Rinc\'on, D. (2026). The Deficiency Indices of $D$. Preprint.

\bibitem{one_wall}
Rinc\'on, D. (2026). One Wall: The Local-to-Global Conjecture. Preprint.

\bibitem{arakelov1974}
Arakelov, S.~J. (1974).
An intersection theory for divisors on an arithmetic surface.
\textit{Izvestiya Akademii Nauk SSSR}, 38(6), 1179--1192.

\bibitem{fargues2021}
Fargues, L. and Fontaine, J.-M. (2021).
\textit{Courbes et fibrés vectoriels en théorie de Hodge $p$-adique}.
Astérisque 406, Société Mathématique de France.

\bibitem{weil1948}
Weil, A. (1948).
\textit{Sur les courbes algébriques et les variétés qui s'en déduisent}.
Hermann, Paris.

\bibitem{deligne1974}
Deligne, P. (1974).
La conjecture de Weil, I.
\textit{Publications Mathématiques de l'IHÉS}, 43, 273--307.

\end{thebibliography}

\end{document}
